\section{Exercises}
\label{sec:2.4}

\begin{enumerate}

\item[2.1] Show that, if $U \sim \Unif(0,1)$, then
\[
  \left\lceil \frac{\log(U)}{\log(1-p)} \right\rceil \sim \Geom(p),
\]
where $\lceil a \rceil$ denotes the ceiling of $a \in \R$, namely the least integer greater than or equal to $a$.
Use this result to transform $N = 10^4$ draws from a uniform distribution into $N$ draws from a
$\Geom(p)$ with $p = 0.2$. Plot a histogram of the transformed sample along with the p.m.f.\ of
a $\Geom(0.2)$ random variable.

\item[2.2] For an arbitrary random variable $X$ with c.d.f.\ $F$, define the generalized inverse of $F$ by
\[
  F^-(u) := \inf\{x : F(x) \ge u\}.
\]
Show that if $U \sim \Unif(0,1)$, then $F^-(U) \overset{d}{=} X$.

\item[2.3] Let $U^{(1)}, U^{(2)} \iid \Unif(0,1)$, and set
\begin{align*}
  X^{(1)} &= \sqrt{-2\log(U^{(1)})}\cos\!\bigl(2\pi U^{(2)}\bigr), \\
  X^{(2)} &= \sqrt{-2\log(U^{(1)})}\sin\!\bigl(2\pi U^{(2)}\bigr).
\end{align*}

\begin{enumerate}[label=\alph*)]
  \item Show that $X^{(1)}, X^{(2)}$ are independent $\mathcal{N}(0,1)$ samples. This is known as the Box-Muller
  transformation method to obtain samples $X^{(1)}, X^{(2)} \iid \mathcal{N}(0,1)$.

  \item Find the KR rearrangement from $\Unif(0,1)^2$ to the standard bivariate normal $\mathcal{N}(0, I_2)$. Compare the running time with the Box-Muller method to obtain $N = 10^5$ samples from the
  standard bivariate normal distribution.
\end{enumerate}

\item[2.4] Prove Theorem~\ref{thm:2.4} and generalize it to characterize the KR rearrangement
between strictly positive p.d.f.\ $g$ and strictly positive p.d.f.\ $f$.

\item[2.5] Rejection sampling can be implemented as long as the ratio $f/g$ can be evaluated
up to a constant, $f/g \propto \tilde{f}/\tilde{g}$ provided that an upper bound $\tilde{M}$ can be found on $\tilde{f}/\tilde{g}$. Spell out
the details on how to do rejection sampling in this case, and explain why the algorithm produces
samples from $f$. Recall that the probability of acceptance in an accept/reject algorithm with
upper bound $M$ on the density ratio $f/g$ is $1/M$. \textbf{Warning}: the acceptance probability for the
algorithm that you have just described using the unnormalized ratio $\tilde{f}/\tilde{g}$ is \textit{not} $1/\tilde{M}$. Show that
the expected value of the acceptance rate, $\mathbb{E}\!\left[\mathbf{1}_{\{U < \tilde{f}/\tilde{M}\tilde{g}\}}\right]$, can be used to compute the missing
constant in $f/g$.

\item[2.6] Recall that a random variable $X \sim \GammaDist(\alpha, \beta)$ with $\alpha > 0$, $\beta > 0$ has
p.d.f.\ $f(x) = \dfrac{\beta^\alpha}{\Gamma(\alpha)} x^{\alpha-1} e^{-\beta x}$. Following Example~\ref{ex:2.8}, use ABC to sample from a posterior
distribution with $\Exp(\lambda)$ likelihood function and $\GammaDist(0.4, 0.4)$ prior. Draw $n_o = 100$
observations from an $\Exp(0.4)$, and define $y_o$ to be their sum (which is sufficient statistic
for $\lambda$). Use the distance function $\dfrac{1}{n_o}|y^{(n)} - y_o|$ (where $y^{(n)}$ is the sum of the 100 observations
simulated from the likelihood in the $n$-th iteration of the algorithm) and three tolerance values
$\varepsilon = 0.005, 0.1, 0.2$. For each tolerance, plot the histogram of the approximate posterior samples
along with the true posterior distribution, which you should find analytically. Further, report the
number of generated samples used for each tolerance in order to obtain $N = 1000$ accepted
samples.

\end{enumerate}
